\appendix
\chapter{System Architecture and Technical Specifications}

This appendix provides detailed technical documentation, configuration specifications, API contracts, and implementation details for the MyAIStorybook system.

\section{JSON Story Structure Schema}

The system employs a standardized JSON schema for story representation throughout the pipeline, validated by Pydantic models defined in \texttt{backend/models/story\_schema.py}.

\subsection{Complete Story Schema Example}

\begin{verbatim}
{
  "title": "The Adventure of the Brave Little Fox",
  "setting": "A magical forest where animals can talk",
  "characters": [
    "Felix the Fox",
    "Olivia the Owl"
  ],
  "scenes": [
    {
      "scene_number": 1,
      "text": "In a magical forest where animals could talk, 
              lived a small orange fox named Felix with 
              bright green eyes. Every morning, Felix would 
              wake up wondering what exciting adventure the 
              day might bring...",
      "image_description": "Small orange fox with green eyes 
                            in forest",
      "image_path": "generated/images/brave_fox_1697234567_
                     scene_1.png"
    },
    {
      "scene_number": 2,
      "text": "One sunny morning, Felix met Olivia, a wise old 
              owl with silver feathers. She told him about 
              a hidden treasure deep in the forest...",
      "image_description": "Fox looking up at silver owl 
                            in tree",
      "image_path": "generated/images/brave_fox_1697234567_
                     scene_2.png"
    },
    {
      "scene_number": 3,
      "text": "With Olivia's guidance and his own courage, 
              Felix discovered that the real treasure was 
              the friendship they built along the journey...",
      "image_description": "Fox and owl standing together 
                            smiling",
      "image_path": "generated/images/brave_fox_1697234567_
                     scene_3.png"
    }
  ],
  "meta": {
    "prompt_type": "normal",
    "enhanced_prompt": "A courageous young fox embarks on 
                        adventure..."
  }
}
\end{verbatim}

\section{API Endpoint Specifications}

\subsection{POST /api/generate}

\textbf{Description:} Initiates story generation with the specified prompt and configuration parameters.

\textbf{Request Body:}
\begin{verbatim}
{
  "prompt": "string (required)",
  "generate_images": boolean (optional, default: true),
  "use_personalized_images": boolean (optional, default: false),
  "user_photo": "base64_encoded_image (optional)",
  "mode": "simple | personalized (optional, default: simple)"
}
\end{verbatim}

\textbf{Response:}
\begin{verbatim}
{
  "story": Story object (JSON schema above),
  "pdf_path": "string (path to generated PDF)",
  "status": "success | error"
}
\end{verbatim}

\subsection{POST /api/chat}

\textbf{Description:} Interactive chat mode for conversational interaction with story characters.

\textbf{Request Body:}
\begin{verbatim}
{
  "message": "string (user chat message)",
  "story_id": "string (required)",
  "session_id": "string (optional)"
}
\end{verbatim}

\textbf{Response:}
\begin{verbatim}
{
  "response": "string (chatbot response)",
  "conversation_id": "string"
}
\end{verbatim}

\subsection{POST /api/workshop/session}

\textbf{Description:} Creates a new idea workshop session for interactive story development.

\textbf{Request Body:}
\begin{verbatim}
{
  "mode": "new_idea | improvement (required)",
  "existing_story": "string (optional, for improvement mode)"
}
\end{verbatim}

\textbf{Response:}
\begin{verbatim}
{
  "session_id": "string",
  "initial_message": "string"
}
\end{verbatim}

\subsection{POST /api/workshop/message}

\textbf{Description:} Sends a message in an active workshop session.

\textbf{Request Body:}
\begin{verbatim}
{
  "session_id": "string (required)",
  "message": "string (required)"
}
\end{verbatim}

\textbf{Response:}
\begin{verbatim}
{
  "response": "string (agent response)",
  "ready_to_generate": boolean,
  "metadata": object (collected story requirements)
}
\end{verbatim}

\subsection{GET /api/stories}

\textbf{Description:} Retrieves list of user's generated stories with metadata.

\textbf{Query Parameters:}
\begin{verbatim}
?limit=10&offset=0&filter=recent
\end{verbatim}

\textbf{Response:}
\begin{verbatim}
{
  "stories": [
    {
      "id": "string",
      "title": "string",
      "created_at": "timestamp",
      "mode": "simple | personalized"
    }
  ],
  "total": integer
}
\end{verbatim}

\subsection{GET /device}

\textbf{Description:} Returns hardware information including device type and GPU availability.

\textbf{Response:}
\begin{verbatim}
{
  "device": "cuda | cpu",
  "name": "string (GPU name if available)"
}
\end{verbatim}

\section{Database Schema}

The system uses PostgreSQL with SQLAlchemy ORM. The database consists of 7 tables:

\subsection{Users Table}
\begin{itemize}
    \item \texttt{id}: Primary key (integer)
    \item \texttt{username}: Unique username (string)
    \item \texttt{email}: User email (string)
    \item \texttt{hashed\_password}: Bcrypt hashed password (string)
    \item \texttt{created\_at}: Account creation timestamp
\end{itemize}

\subsection{Stories Table}
\begin{itemize}
    \item \texttt{id}: Primary key (integer)
    \item \texttt{user\_id}: Foreign key to Users table
    \item \texttt{title}: Story title (string)
    \item \texttt{mode}: Generation mode (simple/personalized)
    \item \texttt{story\_data}: Complete story JSON (JSONB)
    \item \texttt{created\_at}: Story creation timestamp
\end{itemize}

\subsection{Chat Conversations \& Messages Tables}
\begin{itemize}
    \item \texttt{ChatConversations}: Links conversations to stories
    \item \texttt{ChatMessages}: Stores individual chat messages with role (user/assistant)
\end{itemize}

\subsection{Workshop Sessions, Messages \& Stories Tables}
\begin{itemize}
    \item \texttt{WorkshopSessions}: Tracks idea workshop sessions
    \item \texttt{WorkshopMessages}: Stores workshop conversation history
    \item \texttt{WorkshopStories}: Generated stories from workshop sessions
\end{itemize}

\section{stable-diffusion-webui API Integration}

The system integrates with AUTOMATIC1111's stable-diffusion-webui local API for image generation with IP-Adapter support.

\subsection{Key Endpoints Used}

\textbf{/sdapi/v1/txt2img:} Initial image generation from text prompt with IP-Adapter embedding injection.

\textbf{/sdapi/v1/img2img:} Image refinement and upscaling pass for quality enhancement.

\subsection{Configuration Parameters}
\begin{itemize}
    \item Model: Realistic Vision v5.1
    \item Sampler: DPM++ 2M Karras
    \item Steps: 25-30
    \item CFG Scale: 7-7.5
    \item Resolution: 512x512 (txt2img) → 768x768 (img2img)
    \item IP-Adapter: Enabled with user photo embedding
\end{itemize}

\section{Future Work}

The following features are planned for future iterations:

\subsection{Text-to-Speech (TTS) Integration}
\begin{itemize}
    \item Page-by-page voice narration using Coqui TTS
    \item Character-specific voice customization
    \item Audio book export functionality
\end{itemize}

\subsection{Speech-to-Text (STT) Integration}
\begin{itemize}
    \item Voice input processing using OpenAI Whisper
    \item Real-time speech recognition
    \item Multi-language support
\end{itemize}

\subsection{Additional Enhancements}
\begin{itemize}
    \item Multi-language story generation
    \item Advanced character customization controls
    \item Story templates for specific genres
    \item Collaborative storytelling features
    \item Mobile applications (iOS/Android)
    \item Cloud deployment option
    \item Story analytics and recommendation system
\end{itemize}
